\documentclass[11pt]{scrartcl}

\usepackage[stixtwo]{fontsetup}
\usepackage{amsmath}

\begin{document}
$x^2 - x - 6 = (x - 3)(x + 2)$ és $x^2 - 4 = (x - 2)(x + 2)$, így
$x \ne \pm 2$ esetén \(f(x) = \frac{x - 3}{x - 2}\).
\begin{enumerate}
	\item Igazolni kell, hogy minden $\varepsilon > 0$ számhoz létezik
	      $\delta(\varepsilon) > 0$, hogy $0 < |x - 4| < \delta(\varepsilon)$
	      esetén $\left| f(x) - \frac{1}{2} \right| < \varepsilon$. Legyen
	      $\varepsilon > 0$.
	      \[
		      \left| \frac{x - 3}{x - 2} - \frac{1}{2} \right|
		      = \frac{|2x - 6 - x + 2|}{2|x - 2|}
		      = \frac{|x - 4|}{2|x - 2|}
	      \]
	      Ha $|x - 4| < 1$, akkor $3 < x < 5$, így $|x - 2| > 1$, tehát
	      \[
		      \frac{|x - 4|}{2|x - 2|} < \frac{|x - 4|}{2} < \varepsilon,
	      \]
	      ha $|x - 4| < 2\varepsilon$. Így
	      $\delta(\varepsilon) = \min(1, 2\varepsilon)$.

	\item \[
		      \lim_{x \to -2} f(x)
		      = \lim_{x \to -2} \frac{x - 3}{x - 2}
		      = \frac{-5}{-4}
		      = \frac{5}{4}
	      \]

	\item \[
		      \lim_{x \to +\infty} f(x)
		      = \lim_{x \to +\infty} \frac{1 - \frac{3}{x}}{1 - \frac{2}{x}}
		      = 1
	      \]
\end{enumerate}
\end{document}
